\documentclass[10pt,a4paper]{amsart}
\usepackage{amsmath,amssymb,amsthm,graphicx}
\usepackage{verbatim}
\newenvironment{reference}{\comment}{\endcomment}
%\newenvironment{reference}{}{}
\newenvironment{slogan}{\comment}{\endcomment}
\newenvironment{history}{\comment}{\endcomment}

% For commutative diagrams we use Xy-pic
\usepackage[all]{xy}

% We use 2cell for 2-commutative diagrams.
\xyoption{2cell}
% Theorem environments.
%
\theoremstyle{plain}
\newtheorem{theorem}[subsection]{Theorem}
\newtheorem{proposition}[subsection]{Proposition}
\newtheorem{lemma}[subsection]{Lemma}

\theoremstyle{definition}
\newtheorem{definition}[subsection]{Definition}
\newtheorem{example}[subsection]{Example}
\newtheorem{exercise}[subsection]{Exercise}
\newtheorem{situation}[subsection]{Situation}

\theoremstyle{remark}
\newtheorem{remark}[subsection]{Remark}
\newtheorem{remarks}[subsection]{Remarks}

\numberwithin{equation}{subsection}

% Macros
%
\def\lim{\mathop{\mathrm{lim}}\nolimits}
\def\colim{\mathop{\mathrm{colim}}\nolimits}
\def\Spec{\mathop{\mathrm{Spec}}}
\def\Hom{\mathop{\mathrm{Hom}}\nolimits}
\def\Ext{\mathop{\mathrm{Ext}}\nolimits}
\def\SheafHom{\mathop{\mathcal{H}\!\mathit{om}}\nolimits}
\def\SheafExt{\mathop{\mathcal{E}\!\mathit{xt}}\nolimits}
\def\Sch{\mathit{Sch}}
\def\Mor{\mathop{\mathrm{Mor}}\nolimits}
\def\Ob{\mathop{\mathrm{Ob}}\nolimits}
\def\Sh{\mathop{\mathit{Sh}}\nolimits}
\def\NL{\mathop{N\!L}\nolimits}
\def\CH{\mathop{\mathrm{CH}}\nolimits}
\def\proetale{{pro\text{-}\acute{e}tale}}
\def\etale{{\acute{e}tale}}
\def\QCoh{\mathit{QCoh}}
\def\Ker{\mathop{\mathrm{Ker}}}
\def\Im{\mathop{\mathrm{Im}}}
\def\Coker{\mathop{\mathrm{Coker}}}
\def\Coim{\mathop{\mathrm{Coim}}}

% Boxtimes
%
\DeclareMathSymbol{\boxtimes}{\mathbin}{AMSa}{"02}

%
% Macros for moduli stacks/spaces
%
\def\QCohstack{\mathcal{QC}\!\mathit{oh}}
\def\Cohstack{\mathcal{C}\!\mathit{oh}}
\def\Spacesstack{\mathcal{S}\!\mathit{paces}}
\def\Quotfunctor{\mathrm{Quot}}
\def\Hilbfunctor{\mathrm{Hilb}}
\def\Curvesstack{\mathcal{C}\!\mathit{urves}}
\def\Polarizedstack{\mathcal{P}\!\mathit{olarized}}
\def\Complexesstack{\mathcal{C}\!\mathit{omplexes}}
% \Pic is the operator that assigns to X its picard group, usage \Pic(X)
% \Picardstack_{X/B} denotes the Picard stack of X over B
% \Picardfunctor_{X/B} denotes the Picard functor of X over B
\def\Pic{\mathop{\mathrm{Pic}}\nolimits}
\def\Picardstack{\mathcal{P}\!\mathit{ic}}
\def\Picardfunctor{\mathrm{Pic}}
\def\Deformationcategory{\mathcal{D}\!\mathit{ef}}
\usepackage{cleveref}
\crefname{originallemma}{original lemma}{original lemmas}
\Crefname{originallemma}{Original lemma}{Original lemmas}
\makeatletter
\newlabel{formal-spaces-lemma-map-completions-representable}{{0APV}{source}{}{sourcepointer-1}{}}
\newlabel{formal-spaces-lemma-map-completions-representable@cref}{{[originallemma][0][]0APV}{[1][0][]source}{}{}{}}
\makeatother
\crefname{originalsection}{original section}{original sections}
\Crefname{originalsection}{Original section}{Original sections}
\makeatletter
\newlabel{formal-spaces-section-finite-type}{{0AM3}{source}{}{sourcepointer-2}{}}
\newlabel{formal-spaces-section-finite-type@cref}{{[originalsection][0][]0AM3}{[1][0][]source}{}{}{}}
\makeatother
\crefname{originallemma}{original lemma}{original lemmas}
\Crefname{originallemma}{Original lemma}{Original lemmas}
\makeatletter
\newlabel{spaces-morphisms-lemma-base-change-finite-type}{{03XH}{source}{}{sourcepointer-3}{}}
\newlabel{spaces-morphisms-lemma-base-change-finite-type@cref}{{[originallemma][0][]03XH}{[1][0][]source}{}{}{}}
\makeatother
\title{Completion of a finite type morphism is finite type}
\author{The Stacks Project Authors}
\date{stacks-project:restricted.tex:local-snapshot-20261012; GNU FDL 1.2 or later}
\begin{document}
\maketitle
\noindent\textit{Editorial note.} All delivered mathematical text is the author's own, copied from the pinned Stacks Project chapter file restricted.tex; no statement or proof step was rewritten, reordered or removed.\par\medskip
\noindent\textit{Editorial note.} The counted proof is the author's complete argument: the counted statement is the author lemma environment labelled lemma-map-completions-finite-type; the proof environment immediately following it is its own proof; the context span delivers the section situation and the cartesian square that the statement's phrase 'in the situation above' refers to\par\medskip
\noindent\textit{Editorial note.} The label lemma-map-completions-finite-type names the counted statement, and the proof environment immediately following it in the source is delivered verbatim as its proof.\par\medskip
\noindent\textit{Editorial note.} The counted statement is carried by the author's own lemma environment in the chapter source.\par\medskip
\noindent\textit{Allowed standard prerequisites (editorial).}\begin{itemize}
\item the /T completion functor and the cartesian square relating f\_\{/T'\} to f (section-completion-functor, delivered as context)
\item base change preserves local finite type (spaces-morphisms-lemma-base-change-finite-type of Morphisms of Spaces)
\end{itemize}\medskip
\ifdefined\hypertarget\hypertarget{sourcepointer-1}{}\fi
\noindent\textit{Source reference.} Original lemma 0APV: sources/stacks-chow/tags.txt, source lines 13488--13488. Prerequisite pointer invoked at 6069; it names a standard fact of a sibling chapter that the counted argument uses, and every step of the counted proof built on it is retained verbatim in the delivered body.\par\medskip
\ifdefined\hypertarget\hypertarget{sourcepointer-2}{}\fi
\noindent\textit{Source reference.} Original section 0AM3: sources/stacks-chow/tags.txt, source lines 13391--13391. Prerequisite pointer invoked at 6016; it names a standard fact of a sibling chapter that the counted argument uses, and every step of the counted proof built on it is retained verbatim in the delivered body.\par\medskip
\ifdefined\hypertarget\hypertarget{sourcepointer-3}{}\fi
\noindent\textit{Source reference.} Original lemma 03XH: sources/stacks-chow/tags.txt, source lines 5148--5148. Prerequisite pointer invoked at 6022; it names a standard fact of a sibling chapter that the counted argument uses, and every step of the counted proof built on it is retained verbatim in the delivered body.\par\medskip
\section{The completion functor}
\label{section-completion-functor}

\noindent
In this section we consider the following situation.
First we fix a base scheme $S$. All rings, topological rings,
schemes, algebraic spaces, and formal algebraic spaces and
morphisms between these will be over $S$. Next, we fix an
algebraic space $X$ and a closed subset $T \subset |X|$.
We denote $U \subset X$ be the open subspace with $|U| = |X| \setminus T$.
Picture
$$
U \to X \quad |X| = |U| \amalg T
$$
In this situation, given an algebraic space $X'$ over $X$, i.e., an
algebraic space $X'$ endowed with a morphism $f : X' \to X$, then we
denote $T' \subset |X'|$ the inverse image of $T$ and we let $U' \subset X'$
be the open subspace with $|U'| = |X'| \setminus T'$. Picture
$$
U' = f^{-1}U
\quad\quad
\vcenter{
\xymatrix{
U' \ar[d] \ar[r] & X' \ar[d]_f \\
U \ar[r] & X
}
}
\quad\quad
\vcenter{
\xymatrix{
|U'| \ar[r] \ar[d] & |X'| \ar[d]^{|f|} & T' \ar[l] \ar[d] \\
|U| \ar[r] & |X| & T \ar[l]
}
}
\quad\quad
T' = |f|^{-1}T
$$
We will relate properties of $f$ to properties of the induced morphism
$$
f_{/T} : X'_{/T'} \longrightarrow X_{/T}
$$
of formal completions. As indicated in the displayed formula, we will
denote this morphism $f_{/T}$. We have already seen that
$f_{/T}$ is representable by algebraic spaces in
Formal Spaces, Lemma
\ref{formal-spaces-lemma-map-completions-representable}.
In fact, as the proof of that lemma shows, the diagram
$$
\xymatrix{
X'_{/T'} \ar[d]_{f_{/T}} \ar[r] & X' \ar[d]^f \\
X_{/T} \ar[r] & X
}
$$
is cartesian. Please keep this fact in mind whilst reading the lemmas stated
and proved below.

\begin{lemma}
\label{lemma-map-completions-finite-type}
In the situation above. If $f$ is locally of finite type, then
$f_{/T}$ is locally of finite type.
\end{lemma}

\begin{proof}
(Finite type morphisms of formal algebraic spaces are discussed in
Formal Spaces, Section \ref{formal-spaces-section-finite-type}.)
Namely, suppose that $Z \to X$ is a morphism from a scheme into $X$ such
that $|Z|$ maps into $T$. From the cartesian square above we see that
$Z \times_X X'$ is an algebraic space representing
$Z \times_{X_{/T}} X'_{/T'}$. Since $Z \times_X X' \to Z$
is locally of finite type by Morphisms of Spaces, Lemma
\ref{spaces-morphisms-lemma-base-change-finite-type} we conclude.
\end{proof}

\end{document}
